\documentclass[10pt,a4paper]{amsart}
\usepackage[margin=19mm]{geometry}
\usepackage{amsmath,amssymb,mathtools}
\usepackage[scr=rsfs,cal=euler]{mathalfa}
\usepackage{xfrac}
\usepackage[unicode,hidelinks,bookmarks=false]{hyperref}
\newcommand{\ie}{i.e.\ }
\newcommand{\cf}{cf.\ }
\newcommand{\resp}{respectively\ }
\newcommand{\Subsection}{\subsection}
\providecommand{\nobreakdash}{\nobreak\hbox{-}}
\setlength{\emergencystretch}{2em}
\multlinegap0pt
\usepackage{amssymb}
\usepackage{enumerate}
\usepackage{xcolor}
\newtheorem{lemma}{Lemma}
\title{Singular Turing bifurcations and spatial canard solutions \\ in nonlinear reaction-diffusion systems}
\author{Robert Jencks, Tasso J. Kaper, Theodore Vo}
\date{Source: arXiv:2609.08969v1 (CC BY 4.0)}
\begin{document}
\maketitle

\noindent\emph{Source and licence.} Singular Turing bifurcations and spatial canard solutions \\ in nonlinear reaction-diffusion systems. Version arXiv:2609.08969v1 (CC BY 4.0), distributed by arXiv under the Creative Commons Attribution 4.0 International licence, http://creativecommons.org/licenses/by/4.0/, as recorded in the arXiv licence metadata for this exact version and shown on the abstract page https://arxiv.org/abs/2609.08969v1. Full licence notice: \texttt{licenses/arxiv-2609.08969-CC-BY-4.0.txt}.\par\smallskip

\noindent\textit{Editorial scope.} Counted unit: the article's lemma lem-persistenttrue+fauxcanards-RFSNII (source0.tex lines 1696--1700). The article's complete original proof of it is delivered with it (source0.tex lines 1702--1726, one \texttt{proof} environment of 25 lines). No mathematics is written, repaired or re-derived: the statement and the proof are the article's own lines, in the article's own notation, and no retained line is substituted or re-flowed.  Retained uncounted as the source-local context the proof needs: lines 279--286; lines 1175--1180; lines 1656--1664; lines 1677--1685.  Nothing was omitted from any retained span, so the package carries no omission record.  No retained line contains a citation, so the wrapper carries no bibliography and no bibliography entry.  No retained line contains a figure environment, an image-inclusion command or drawing code, so no image is delivered and the package declares no asset dependency.  Editorial formatting only, in addition: an AMS \texttt{amsart} preamble that restates the article's own declarations and macros used by the retained lines, namely the environments \texttt{lemma} on separate counters, together with the standard packages those lines need.  The retained environments print in this document as Lemma 1 in order of appearance, whereas the article prints the same items with the numbering of its own full document.  The complete native source of the article as distributed in the arXiv e-print for this exact version is bundled unchanged as \texttt{sources/209764/source0.tex}.
\begin{equation}  \label{eq:y-ODE}
\begin{split}
    u_y &= p \\
    p_y &= v + u^2 + \alpha_1 uv + \alpha_2 u^3 + \mathcal{O}(v^2, u^2v, u^4),  \\
    v_y &= \delta q \\
    q_y &= \delta \left( -\frac{\lambda}{2} + \alpha_3 u + \alpha_4 v + \alpha_5 u^2 + \mathcal{O}(uv,v^2,u^3)\right).
\end{split}
\end{equation}

\begin{lemma} \label{lem-persistence-true+fauxcanards-RFS}
There exists a $\delta_0>0$ sufficiently small such that for each $0<\delta<\delta_0$ and any $\mathcal{O}(1)$ value of $\lambda>0$,
the singular true and faux canards $\Gamma_t^0$ and $\Gamma_f^0$ persist as maximal true and faux canards, labeled $\Gamma_t^\delta$ and $\Gamma_f^\delta$, of \eqref{eq:y-ODE}.
The maximal canards lie in the transverse intersection of slow invariant manifolds,
and the connection between $\Gamma_t^\delta$ and $\Gamma_f^\delta$ persists for $0 < \delta< \delta_0$.
\end{lemma}

\begin{equation}\label{eq:limitGamma20}
   \begin{split}
       \kappa_{21}(\Gamma_{20}(\hat{y}_2)) = \left\{ 0, 
       -\frac{4\sqrt{3}\alpha_3 \hat{y}_2}{ z_2^{3/2}}, 
       -1-\frac{\alpha_3}{24 z_2^2}, 
       \frac{-2 \alpha_3^2\hat{y}_2^3}{\sqrt{3}z_2^{3/2}}, 
       \frac{12}{z_2}, \frac{24 \sqrt{3} \hat \lambda_2}{z_2^{3/2}} \right\}, 
   \end{split} 
\end{equation}

\begin{equation}\label{eq:limit-tanvectors}
\lim_{\hat{y}_2 \to \pm \infty} 
\frac{\frac{d}{d\hat{y}_2}\kappa_{21}(\Gamma_{20})}{\Vert \frac{d}{d\hat{y}_2}\kappa_{21}(\Gamma_{20})\Vert}
= \left(0, 
\frac{\alpha_3}{\sqrt{-\alpha_3(3-\alpha_3)}}, 
0, 0,
\mp \sqrt{\frac{3}{3 - \alpha_3}},  
0 \right) \in \mathbb E^c(\ell_1).
\end{equation}

\begin{lemma}\label{lem-persistenttrue+fauxcanards-RFSNII}
There exists a $\delta_0>0$ sufficiently small such that for each $0< \delta<\delta_0$ and each $\lambda\ge 0$ sufficiently small,
the singular true canard $\Gamma_{t}^0$ and singular faux canard $\Gamma_{f}^0$ of $M_{\rm RFSN-II}$ perturb to maximal true and faux canard solutions,  which we label $\Gamma_{t}^{\delta}$ and $\Gamma_{f}^{\delta}$ respectively, of \eqref{eq:y-ODE}.
Also, for positive values of $\lambda > \mathcal{O}(\delta)$, the connection between $\Gamma_t^\delta$ and $\Gamma_f^\delta$ persists for $0<\delta<\delta_0$.
\end{lemma}

\begin{proof} 
In chart $\hat{K}_2$, the left-half ($-\infty < \hat{y}_2<0$) of the key algebraic solution $\Gamma_{20}(\hat{y}_2)$ corresponds to the singular true canard $\Gamma_{t}^0$ and the right-half ($0< \hat{y}_2 < \infty$) to the singular faux canard $\Gamma_{f}^0$.
By \eqref{eq:limitGamma20},
$\Gamma_{20}(\hat{y}_2)$ is backward asymptotic to the equilibrium ${\bf p}_{-} \in \ell_1 \subset \mathcal{M}_1$, and forward asymptotic to ${\bf p}_{+} \in \ell_1 \subset \mathcal{M}_1$.
Both points lie on the equator of the blown-up hemisphere $\mathbb{S}^5$.

The analysis of the asymptotics of the tangent vector in \eqref{eq:limit-tanvectors} shows that $\Gamma_{20}(\hat{y}_2)$ emanates from the four-dimensional normally attracting $W^{c}\left( \bf{p}_- \right)$ and terminates in the four-dimensional normally repelling $W^{c}\left( \bf{p}_+ \right)$ in the $\left\{ \hat{r}_1=0 \right\}$ hyperplane of chart $\hat{K}_1$.
The asymptotic behavior is decomposed into components tangent to 
$\mathcal{M}_{1}\rvert_{\hat{r}_{1}=0}$ and in the neutrally stable fiber over the limit points.
Hence, the only solution of the variational equation about $\Gamma_{20}$ that is bounded in both limits as $\hat{y}_2 \to \pm \infty$ is $\tfrac{d}{d\hat{y}_2}\kappa_{21}(\Gamma_{20})$.
All other solutions diverge from ${\bf p}_-$ in backward time and from ${\bf p}_+$ in forward time. 
Therefore, in the invariant set $\{ \hat r_2=0 \}$, $\Gamma_{20}(\hat{y}_2)$ lies in the transverse intersection of $W^{c}\left( \bf{p}_- \right)$ and $W^{c}\left( \bf{p}_+ \right)$, which is a two-dimensional intersection manifold.

Next, because this intersection is transverse on $\{ \hat r_2 = 0 \}$ for each fixed value of $\hat \lambda_2$, we know from regular perturbation theory that the intersection of $W^{c}({\bf p}_-)$ and $W^{c}({\bf p}_+)$ is also transverse in chart $\hat K_2$ for each $\hat r_2>0$ sufficiently small. 
By tracking this intersection into chart $\hat K_1$, we see that the intersection remains transverse.
Moreover, the intersection manifold $W^{c}\left( \bf{p}_- \right) \cap W^{c}\left( \bf{p}_+\right)$ is foliated by the parameter $\mathcal{\hat{\lambda}}_2$ in chart $\hat{K}_2$.
Hence, by tracking the solutions of this foliation into chart $\hat{K}_1$, we see that for each fixed $\hat{r}_2$ 
this foliation carries over into a nontrivial foliation by $\hat{r}_1^2 {\hat{\lambda}}_1 = {\rm constant}$ in $\hat{K}_1$
(where we recall that ${\hat{\lambda}}_1 \hat{r}_1^2 = {\hat{\lambda}}_2 \hat{r}_2^2$).
Therefore, for each fixed $\mathcal{\hat{\lambda}}_2$ (or $\hat{r}_1^2 \mathcal{\hat{\lambda}}_1$ with $\hat{r}_1$ sufficiently small), $\Gamma_{20}(\hat{y}_2)$ is a curve, and it corresponds to the maximal canard solution. 
This shows that the singular canards $\Gamma_{t}^0$ and $\Gamma_{f}^0$ perturb to maximal canard solutions $\Gamma_{t}^{\delta}$ and $\Gamma_{f}^{\delta}$ for all $\delta > 0$ sufficiently small, as stated in the lemma.

To establish the persistence of the connection, we recall that the last part of the proof of Lemma~\ref{lem-persistence-true+fauxcanards-RFS} establishes the persistence of the connection between the maximal true and faux canards in the regime of $\mathcal O(1)$ $\lambda>0$, where the folded singularity is an RFS.
The analysis performed there may be extended down to the regime in which $\lambda=\mathcal{O}(\delta)$ by using a $\delta$-dependent rescaling of $q$.
\end{proof}

\end{document}
